\begin{titlepage}
\raggedright
{\color{teal}\sffamily\bfseries THE ENERGY BALANCE PROJECT\par}
\vspace{30pt}
{\color{navy}\sffamily\bfseries\fontsize{26}{31}\selectfont Mathematics, Implementation and Evidence\par}
\vspace{20pt}
{\Large Part II\par}
\vspace{12pt}
{\large An integrated explanatory edition\par}
\vspace{20pt}
{\large Konrad Grzyb\par}
\vfill
{\color{slate}From an exact physical account to a scientific claim that can be checked.\par}
\vspace{16pt}
Author-review manuscript. New Gaussian exposition and explicitly scoped
historical readings. No new scientific experiment is reported by this revision.
\end{titlepage}

\chapter*{Author, rights and review status}
Konrad Grzyb is the author and originator of the Energy Balance Project.
This manuscript was prepared with AI assistance under his direction.
It is supplied for author review, not as a claim of author acceptance,
independent peer review or publication.

The software repository has a separate MIT licence. That software licence
does not automatically grant rights in these books. This revision introduces
no new rights grant. The original editions remain preserved; their results
remain statements about their declared models.

All new numerical examples are illustrative mathematics on explicitly
supplied states, not measured data, operating recommendations, registered
parameters or scientific trajectories. Historical readings retain their
own experimental status; no new study is executed by reproducing them.

\chapter*{Why these two books now belong together}
Part I introduced a question: can an account of physical action describe
changes in declared functional deviation without confusing those changes
with prices, ownership or a promise of recovery? It developed the objects
needed to ask that question and finished with exact finite EBU. This volume
starts where that introduction ends. Its task is to show how an equation
becomes a precise claim, how software can preserve or corrupt that claim,
and how evidence can support only some of the conclusions a reader might
be tempted to draw.

The former Parts II and III separated mathematical structure from code and
evidence. That separation made sense as a classification of subjects, but it
also interrupted the argument. Readers encountered a guarantee in one book
and the implementation assumption on which it depended in another. The
combined volume repairs that connection. Mathematical definitions come first;
group and path accounting precede any discussion of persistent capacity;
physical permission precedes execution order; tests precede interpretation.
The historical experiments then become a worked lesson in asking questions,
not a collection of success claims placed after the mathematics.

This is not a short digest of the two originals. The new explanatory spine
is joined to complete, identified historical readings at the point where
they help the argument. In particular, the detailed dissipation and numerical
proofs, preservation and route reasoning, software safeguards, experimental
sequence, flat outcomes and reproduction discipline remain available.
Old assumptions are kept visible because concealing them would make the
record less intelligible, not more modern.

\chapter*{One continuous course, two kinds of reading}
Main chapters continue Part I's numbering: 33 through 62. These chapters use
the prospective local Gaussian domain unless a passage declares another
model. References and scales are fixed; the first physical domain is
lossless and conserves one homogeneous scalar. A controller is not smuggled
into the definition of the potential.

Historical readings are marked H44, H45 and so on, after the current chapter
that contains them. Each page carries a historical banner and identifies its
original part and PDF page. Its body, equations and internal numbering are
preserved. A reference to original Equation 38.4 inside a reading is therefore
an original-source reference, not a reference to the new Chapter 38. New
running page numbers belong to this integrated volume. Original cross-links
are not copied into destinations where they would become misleading.

Read a historical proof with a pencil beside its assumptions. Some steps,
such as endpoint cancellation, apply to many potential families. Other steps
require a specified response law, a regenerative source classification, a
particular numerical certificate or a frozen experiment. The surrounding
current text explains those boundaries. ``Historical'' does not mean false.
It means that the object and authority of the claim remain identifiable.

The supplied original PDFs remain unchanged. Their editable manuscripts were
not recovered. New chapters are editable source; historical readings retain
their original typesetting with smaller displayed chapter titles and new
margin furniture. The source map records what was reconstructed, what was
retained as a reading, and what remains only in the original. In particular,
the old threshold laboratory is replaced by a Gaussian worked laboratory;
its numerical outcomes are not converted into Gaussian findings.

\chapter*{The dependency route}
Chapters 33--39 establish the domain, units, geometry, physical graph, local
support, force and finite quote. The distinction between force and flux is
explained here, including why Onsager's modelling pattern was useful and
why borrowing it is not a thermodynamic proof about an economy.

Chapters 40--43 compose actions without double counting. A group value is
calculated before a child attribution; attribution is defined before an
owner account; external changes are separated before an invariant is written.
This order prevents a correct sum identity from appearing to prove an
unannounced settlement rule.

Chapters 44--46 examine the additional assumptions needed for dynamics,
permission, routes and long-run reasoning, then work through a detailed
Gaussian laboratory. The laboratory is fixed-state arithmetic, not an
executed scientific trajectory or a chosen campaign configuration.

Chapters 47--56 translate those distinctions into software responsibilities:
schemas, permission, chronology, menus, evaluation, accounts, randomness,
tests and provenance. Chapters 57--60 recover the historical engine and
evidence in enough detail to understand both informative results and limits.
Chapters 61--62 end with what a future study could establish and an audit
that makes missing authority visible.

\begin{lawbox}{Four statements that must never collapse into one}
A conditional identity is a mathematical result. A faithful implementation
is a software achievement. A registered design is a plan. An observed
scientific outcome requires an actual authorized run and its preserved
evidence. Passing from one to the next needs additional work; a book's
chapter order is not that authorization.
\end{lawbox}

\chapter*{How this volume changes the series}
The current series has eight parts, not nine. Former Parts II and III become
this Part II. Former Parts IV--IX shift down one number to Parts III--VIII.
Part I keeps its role and chapter sequence. Old editions, citations, commit
records and historical chapter names retain their original identifiers;
renumbering a forward plan is not rewriting a scientific source.

Only the first two parts are being delivered in this revision. The later
six remain a plan whose subjects depend on the claims they need, not on an
assumption that every earlier experiment will succeed. The closing series
map gives both the new numbering and its old equivalent.
